All maps are linear in $a$ and $X$, so it suffices to take $a = e_i$ and $X = e_K$, with $e_1, \dots, e_n$ an orthogonal basis (\claimref{std-algebra.grade-decomposition}) and $|K| = k$. By \claimref{std-algebra.orthogonal-anticommute}, $e_i$ anticommutes with each factor of $e_K$ other than itself. If $i \notin K$, then $e_i e_K = \pm e_{K \cup \{i\}}$ has grade $k + 1$, and $e_K e_i = (-1)^{k} e_i e_K$. If $i \in K$, then $e_i e_K$ is a multiple of $e_{K \setminus \{i\}}$, of grade $k - 1$, and $e_K e_i = (-1)^{k-1} e_i e_K$, because $e_i$ passes the $k - 1$ other factors and commutes with itself. In both cases $e_i e_K$ is homogeneous, of grade $k + 1$ or $k - 1$, so it equals $e_i \wedge e_K$ or $e_i \rfloor e_K$ respectively and the other part is zero; this gives $aX = a \rfloor X + a \wedge X$. Since $\hat{e}_K = (-1)^{k} e_K$ (\claimref{std-algebra.involution-properties}), $\hat{e}_K e_i = e_i e_K$ when $i \notin K$ and $\hat{e}_K e_i = -e_i e_K$ when $i \in K$, which gives the two symmetrisation formulas. For $X = b$ a vector, $a \rfloor b = \langle ab \rangle_0 = B(a, b)$.
